\section{Absolute values and finite local degrees}

The absolute values in this chapter are multiplicative, real-valued and
satisfy the triangle inequality. The trivial absolute value is allowed.
For an absolute value $v$ on a field $K$, write $K_v$ for its completion.
All statements about sigma algebras refer to the specified sigma algebra;
no countability assumption on the set of places is imposed.
These facts supply the local-field ingredients used in
\cite[\S\S3.3--3.4]{chen-moriwaki}.

\begin{definition}[Adelic structures and properness]\label{def:adelic-structure}
\source{chen-moriwaki-adelic:page-0194,chen-moriwaki-adelic:page-0201}
\lean{ArakelovGeometry.AdelicStructure,ArakelovGeometry.AdelicStructure.instMeasurableSpaceΩ,ArakelovGeometry.AdelicCurve,ArakelovGeometry.AdelicCurve.instFieldK,ArakelovGeometry.AdelicCurve.Ω,ArakelovGeometry.AdelicCurve.v,ArakelovGeometry.AdelicCurve.ϕ,ArakelovGeometry.AdelicCurve.IsProper,ArakelovGeometry.AdelicStructure.toAdelicCurve}
An adelic structure on $K$ consists of a measure space
$(\Omega,\mathcal A,\nu)$ and absolute values $\phi_\omega$ such that
$\omega\mapsto\log\phi_\omega(a)$ is measurable and integrable for every
$a\in K^\times$. An adelic curve is a field together with an adelic
structure. It is proper if
$\int_\Omega\log\phi_\omega(a)\,d\nu(\omega)=0$ for every $a\in K^\times$.
\end{definition}

\begin{definition}[Restriction of absolute values]\label{def:absolute-value-restriction}
\source{chen-moriwaki-adelic:page-0194}
\lean{AbsoluteValue.pullback,AbsoluteValue.restriction}
The pullback of $w$ along an embedding $j:K\to L$ is $a\mapsto w(j(a))$.
In a field extension it is denoted by $w|_K$.
\end{definition}

\begin{lemma}[Maps between completions]\label{lem:completion-embedding}
\uses{def:absolute-value-restriction,lem:completion-nonarchimedean}
\lean{AbsoluteValue.isometry_algebraMap_withAbs,AbsoluteValue.algebraCompletion,AbsoluteValue.algebraMap_completion_def,AbsoluteValue.isometry_algebraMap_completion,AbsoluteValue.norm_algebraMap_completion,AbsoluteValue.coe_smul_completion}
An inclusion $(K,v)\subseteq(L,w)$ with $w|_K=v$ is isometric and extends
to an isometric field embedding $K_v\to L_w$. Multiplication by a scalar
and the embeddings of the original fields commute with these completion maps.
\end{lemma}
\begin{proof}
For $a,b\in K$, $w(a-b)=v(a-b)$, so Cauchy sequences and equivalent
Cauchy sequences remain so under the inclusion. Send the class of $(a_n)$
to its limit in $L_w$. Equality of distances passes to limits and proves
isometry. Addition and multiplication commute with limits, so the extension
is a field homomorphism. It agrees with the original inclusion on constant
sequences. Applying this observation to products of approximating sequences
gives the asserted compatibility with scalar multiplication.
\end{proof}

\begin{lemma}[Completeness and the nonarchimedean property]
\label{lem:completion-nonarchimedean}
\lean{CauSeq.isComplete_completion,AbsoluteValue.isNonarchimedean_completion_iff}
The completion $K_v$ is complete for its induced absolute value $\widehat v$.
Moreover, $v$ is nonarchimedean if and only if $\widehat v$ is
nonarchimedean.
\end{lemma}
\begin{proof}
To verify completeness directly, let $(z_n)$ be a Cauchy sequence in $K_v$.
Choose $a_n\in K$ with $\widehat v(z_n-a_n)<2^{-n}$. The triangle
inequality shows that $(a_n)$ is Cauchy. Its class in the completion is a
limit of $(a_n)$ and hence of $(z_n)$. If $v$ is nonarchimedean, approximate
$x,y\in K_v$ by sequences $a_n,b_n\in K$. Passing to the limit in
$v(a_n+b_n)\leq\max(v(a_n),v(b_n))$ gives the nonarchimedean inequality
for $x,y$. The converse follows by restriction along the isometric
embedding of $K$.
\end{proof}

\begin{lemma}[Archimedean type under extension]\label{lem:archimedean-extension}
\source{chen-moriwaki-adelic:page-0202,chen-moriwaki-adelic:page-0203,chen-moriwaki-adelic:page-0204}
\uses{lem:bounded-integers,lem:completion-nonarchimedean,thm:complete-archimedean,lem:archimedean-rational-normalization}
\notready
For a finite extension $L/K$ and an extending absolute value $w$ of $v$,
$w$ is archimedean if and only if $v$ is archimedean. In that case their
exponents on $\mathbb Q$ agree, and if that exponent is $\kappa$, then
$w^{1/\kappa}$ is an absolute value on $L$.
\end{lemma}
\begin{proof}
A nonarchimedean absolute value restricts to a nonarchimedean one.
Conversely, if $v$ is nonarchimedean, then for every natural integer $n$
one has $w(n)=v(n)\leq1$. Lemma~\ref{lem:bounded-integers} shows that
$w$ is nonarchimedean. In the archimedean case its completion is
archimedean by Lemma~\ref{lem:completion-nonarchimedean}, so the complete
archimedean classification writes
$w=|\sigma(\cdot)|^t$ on $L$, where $\sigma$ takes values in
$\mathbb R$ or $\mathbb C$. On the rational element $2$, restriction
gives $2^t=2^\kappa$, hence $t=\kappa$. It follows that
$w^{1/\kappa}=|\sigma(\cdot)|$, which satisfies the triangle inequality.
\end{proof}

\begin{lemma}[Bounded integers and the ultrametric inequality]
\label{lem:bounded-integers}\notready
An absolute value $v$ is nonarchimedean if and only if $v(n)\leq1$ for every natural
integer $n$. Every absolute value extending a trivial absolute value to an
algebraic extension is trivial.
\end{lemma}
\begin{proof}
If $v$ is nonarchimedean, induction using $n+1=n+1_K$ and $v(1)=1$
gives $v(n)\leq1$ for every natural integer $n$.
Conversely, put $M=\max(v(x),v(y))$. The binomial formula gives
$v(x+y)^n\leq(n+1)M^n$ for every positive integer $n$, since each
binomial coefficient has absolute value at most one. If $M=0$ both $x,y$
vanish. Otherwise take $n$th roots and let $n$ tend to infinity, obtaining
$v(x+y)\leq M$. An extension $w$ of a trivial absolute value
has $w(n)\leq1$ and is therefore nonarchimedean. If a nonzero algebraic
element $a$ had $w(a)>1$, its monic minimal polynomial would give
$w(a)^d\leq\max_{i<d} w(a)^i<w(a)^d$, a contradiction. Thus $w(a)\leq1$;
applying this to $a^{-1}$ gives $w(a)\geq1$.
\end{proof}

\begin{lemma}[Finite-dimensional nontrivial completions]
\label{lem:completion-finite-nontrivial}
\uses{lem:completion-embedding,lem:completion-nonarchimedean}
\lean{AbsoluteValue.finiteDimensional_completion_of_isNontrivial}
If $L/K$ is finite and $v=w|_K$ is nontrivial, then $L_w$ is finite
dimensional over $K_v$.
\end{lemma}
\begin{proof}
Choose a $K$-basis $b_1,\ldots,b_d$ of $L$ and let $W$ be its
$K_v$-linear span in $L_w$. The span is finite dimensional. A finite
dimensional subspace of a normed space over a complete nontrivially normed
field is closed. Thus $W$ is closed in $L_w$. It contains $L$, which is
dense in its completion, so $W=L_w$.
\end{proof}

\begin{lemma}[Generation of finite completions]\label{lem:completion-finite}
\source{chen-moriwaki-adelic:page-0205}
\uses{lem:completion-finite-nontrivial,lem:bounded-integers}
\notready
For every finite extension $L/K$ and every extension $w$ of $v$,
$L_w=K_v(L)$ inside $L_w$, and $[L_w:K_v]\leq[L:K]$.
\end{lemma}
\begin{proof}
For nontrivial $v$, the preceding proof shows that any $K$-basis of $L$
spans $L_w$ over $K_v$. This span is contained in $K_v(L)$ and proves
both claims. If $v$ is trivial, so is $w$ by
Lemma~\ref{lem:bounded-integers}. Both fields are discrete and complete:
a Cauchy sequence is eventually constant, since unequal points have distance
one. Thus $L_w=L$, $K_v=K$, and both claims follow directly.
\end{proof}

\begin{definition}[Local degree]\label{def:local-degree}
\source{chen-moriwaki-adelic:page-0194}
\uses{lem:completion-embedding,lem:completion-finite}
\lean{AbsoluteValue.localDegree}
For a finite extension $L/K$ and an absolute value $w$ on $L$, the local
degree is $n(w/v)=[L_w:K_v]$, where $v=w|_K$ and the scalar embedding is
the completion of $K\hookrightarrow L$.
\end{definition}

\begin{theorem}[Classification of complete archimedean fields]
\label{thm:complete-archimedean}
\source{chen-moriwaki-adelic:page-0203,chen-moriwaki-adelic:page-0204}
\uses{lem:bounded-integers}
\lean{AbsoluteValue.exists_ringEquiv_real_or_complex_of_complete}
A field $F$ complete for an archimedean absolute value $v$ is isomorphic to
$\mathbb R$ or $\mathbb C$, with absolute value the $s$th power of the
usual modulus for some $s\in(0,1]$.
This is the classification invoked in the archimedean part of the proof
of Theorem 3.3.7 of \cite{chen-moriwaki}. The argument below uses the
rational Ostrowski theorem and the real Gelfand--Mazur theorem in
\cite{mathlib-archimedean-classification}.
\end{theorem}
\begin{proof}
Some integer has absolute value greater than one by
Lemma~\ref{lem:bounded-integers}. The characteristic is zero: in positive
characteristic each nonzero element of the prime field has finite
multiplicative order and therefore absolute value one. Ostrowski's theorem
on $\mathbb Q$ gives $v(q)=|q|^s$, $s>0$; the nonarchimedean possibilities
are excluded by that integer. The inequality $v(2)\leq2$ gives $s\leq1$.

The function $u=v^{1/s}$ is an absolute value. Multiplicativity and
definiteness follow from those of $v$. To prove its triangle inequality,
write $A=v(x)^{1/s}$, $B=v(y)^{1/s}$. Binomial expansion gives, for every
positive $n$,
\[
 v(x+y)^n\leq\sum_{i=0}^n
  \bigl({n\choose i}A^iB^{n-i}\bigr)^s
 \leq(n+1)(A+B)^{ns}.
\]
Taking $n$th roots and passing to the limit gives $v(x+y)\leq(A+B)^s$.
The norm $u$ is complete, since raising distances to a positive power
preserves Cauchy sequences and limits. The usual isometric embedding of
$\mathbb Q$ therefore extends to an isometric field embedding
$\mathbb R\to F$. The real Gelfand--Mazur theorem for a normed division
algebra identifies $F$ as an $\mathbb R$-algebra with $\mathbb R$ or
$\mathbb C$. In the real case the scalar embedding is already isometric.
In the complex case let $\psi:\mathbb C\to F$ be the inverse algebra
isomorphism. The identity
$\psi(z)=\operatorname{Re}(z)+\operatorname{Im}(z)\psi(i)$ proves that
$\psi$ is continuous. Thus $u(\psi(z))$ is bounded, say by $B$, on the
unit circle. For $|z|=1$ and every positive integer $n$,
$u(\psi(z))^n\leq B$ and $u(\psi(z))^{-n}\leq B$. Taking $n$th roots
and letting $n$ tend to infinity gives $u(\psi(z))=1$.
Writing a nonzero complex number as its modulus times a unit complex number
then proves that the algebra isomorphism preserves $u$. Finally $v=u^s$.
\end{proof}

\begin{lemma}[Rational normalization at an archimedean place]
\label{lem:archimedean-rational-normalization}
\source{chen-moriwaki-adelic:page-0203,chen-moriwaki-adelic:page-0204}
\uses{lem:completion-nonarchimedean,thm:complete-archimedean}
\lean{AbsoluteValue.charZero_of_not_isNonarchimedean,AbsoluteValue.exists_rpow_eq_of_not_isNonarchimedean}
If $v$ is an archimedean absolute value on $K$, then $K$ has
characteristic zero, and there exists $s\in(0,1]$ such that
$v(q)=|q|^s$ for every $q\in\mathbb Q$.
\end{lemma}
\begin{proof}
The completion remains archimedean, hence is isomorphic to $\mathbb R$
or $\mathbb C$ with absolute value the $s$th power of the usual modulus
for some $s\in(0,1]$. Its subfield $K$ therefore has characteristic zero.
Every field homomorphism between fields of characteristic zero fixes the
prime field $\mathbb Q$. Restricting the norm formula to this prime field
gives the assertion.
\end{proof}

\begin{lemma}[Extension over a complete field]\label{lem:complete-extension}
\source{chen-moriwaki-adelic:page-0208}
\uses{lem:bounded-integers,lem:completion-nonarchimedean,thm:complete-archimedean}
\notready
If $(F,v)$ is complete and $E/F$ is algebraic, there is exactly one
absolute value on $E$ extending $v$.
\end{lemma}
\begin{proof}
For trivial $v$, existence is given by the trivial absolute value and
uniqueness by Lemma~\ref{lem:bounded-integers}. For nontrivial
nonarchimedean $v$, use the spectral norm theorem: on an algebraic extension
of a complete nontrivially normed nonarchimedean field the spectral norm is
multiplicative, extends the base norm, and is the unique extending
multiplicative norm. In the notation $P_a(T)=T^d+\sum_{i<d}c_iT^i$ for
the minimal polynomial, its value is
$\max_{i<d}v(c_i)^{1/(d-i)}$. The extension and uniqueness assertions are
precisely those of the spectral norm theorem, applicable because the base is
complete, its norm is nontrivial, and every $a\in E$ is algebraic. We use
this existing norm theorem in the form recorded in
\cite{mathlib-spectral-norm}.

In the archimedean case, the classification of complete archimedean
absolute-valued fields identifies $F$ with $\mathbb R$ or $\mathbb C$ and
$v$ with the $s$th power of the usual modulus, $0<s\leq1$.
An algebraic extension of $\mathbb C$ is itself, and an algebraic extension
of $\mathbb R$ is $\mathbb R$ or $\mathbb C$; extending the modulus to
$\mathbb C$ and raising it to $s$ gives existence. To check uniqueness,
complete $E$ with respect to any extending absolute value. Its completion
is again archimedean, so the same classification identifies it with
$\mathbb R$ or $\mathbb C$ with exponent $t$. Evaluation on $2$ gives
$t=s$. The induced embedding of $\mathbb R$ fixes the rationals and is
continuous, hence fixes $\mathbb R$. If $E=\mathbb C$, the image of $i$
is $i$ or $-i$, and either choice preserves the usual modulus. Thus every
extending absolute value is the one constructed above.
\end{proof}

\begin{proposition}[Local factorization]\label{lem:local-factorization}
\source{chen-moriwaki-adelic:page-0194,chen-moriwaki-adelic:page-0202,chen-moriwaki-adelic:page-0205}
\uses{lem:complete-extension,lem:completion-finite,def:local-degree}
\notready
Let $L/K$ be finite separable, let $v$ be an absolute value on $K$, and
let $X_v$ be the set of extending absolute values on $L$. Then $X_v$ is
finite and nonempty, and
\[
 K_v\otimes_KL\simeq\prod_{w\in X_v}L_w,\qquad
 \sum_{w\in X_v}[L_w:K_v]=[L:K].
\]
For every $b\in L^\times$,
\[
 v(N_{L/K}(b))=\prod_{w\in X_v}w(b)^{[L_w:K_v]}.
\]
\end{proposition}
\begin{proof}
Take a primitive element $a$ and its separable monic minimal polynomial
$P\in K[T]$. Write $P=\prod_{j=1}^rP_j$ as its distinct monic irreducible
factors over $K_v$. Each field $E_j=K_v[T]/(P_j)$ has the unique extending
absolute value of Lemma~\ref{lem:complete-extension} and is complete,
being finite dimensional over a complete field (in the trivial case it is
discrete). The embedding $L=K[T]/(P)\to E_j$ has dense image: the powers
of the residue of $T$ span $E_j$ and each coefficient in $K_v$ can be
approximated by coefficients in $K$. Its completion is consequently $E_j$.

Conversely, for an extending $w$, Lemma~\ref{lem:completion-finite}
gives $L_w=K_v(a)$. The evaluation homomorphism $K_v[T]\to L_w$ is
surjective and its kernel is generated by one of the $P_j$. Uniqueness of
the extending absolute value identifies $w$ with the value just obtained
from that factor. If two factors gave the same $w$, their maps to the
completion of the same valued field $L$ would agree on the dense subfield
$L$ and on $K_v$, hence would have the same kernel. Thus their factors
would agree. The factors therefore correspond bijectively to $X_v$.

The Chinese remainder theorem gives
$K_v\otimes_K L\simeq K_v[T]/(P)\simeq\prod_jE_j$.
Taking dimensions proves the degree sum. Determinants of multiplication
commute with extension of scalars and with products, so
$N_{L/K}(b)=\prod_jN_{E_j/K_v}(b)$ in $K_v$.
The norm in each separable $E_j/K_v$ is the product over its embeddings
into an algebraic closure of $K_v$. All embeddings preserve the unique
extending absolute value. Taking absolute values therefore gives
$v(N_{E_j/K_v}(b))=w_j(b)^{[E_j:K_v]}$ and proves the last formula.
\end{proof}

\begin{lemma}[Number of extending absolute values]\label{lem:finite-extension-count}
\source{chen-moriwaki-adelic:page-0194,chen-moriwaki-adelic:page-0204}
\uses{lem:local-factorization,lem:finite-separable-part,lem:inseparable-extension-unique}
\notready
For any finite extension $L/K$, an absolute value of $K$ has at least one
and at most $[L:K]$ extensions to $L$.
\end{lemma}
\begin{proof}
Let $E$ be the separable closure of $K$ in $L$. Restriction gives a
bijection between extensions to $L$ and extensions to $E$, because $L/E$
is purely inseparable and the extension formula is unique. The latter
set is nonempty by local factorization. Its cardinality is at most
$[E:K]$, since the positive local degrees sum to $[E:K]$. Finally
$[E:K]\leq[L:K]$ by the dimension formula in the tower.
\end{proof}

\begin{lemma}[Archimedean exponent]\label{lem:archimedean-bound}
\source{chen-moriwaki-adelic:page-0202,chen-moriwaki-adelic:page-0203,chen-moriwaki-adelic:page-0204}
\uses{def:adelic-structure,lem:archimedean-rational-normalization}
\lean{ArakelovGeometry.AdelicStructure.Ωfin,ArakelovGeometry.AdelicStructure.Ωinf,ArakelovGeometry.AdelicStructure.κ,ArakelovGeometry.AdelicStructure.κ_mem_Icc,ArakelovGeometry.AdelicStructure.κ_eq_zero_iff,ArakelovGeometry.AdelicStructure.κ_def,ArakelovGeometry.AdelicStructure.measurableSet_Ωinf,ArakelovGeometry.AdelicStructure.measurable_κ,ArakelovGeometry.AdelicStructure.integrable_κ,ArakelovGeometry.AdelicCurve.κ,ArakelovGeometry.AdelicCurve.Ωinf,ArakelovGeometry.AdelicCurve.Ωfin}
The archimedean locus is measurable. Let $\kappa(\omega)$ be zero on
the nonarchimedean locus and the exponent of the usual absolute value on
$\mathbb Q$ at an archimedean place. Then $0\leq\kappa\leq1$,
$\kappa$ vanishes exactly at the nonarchimedean places, and $\kappa$ is
measurable and integrable.
\end{lemma}
\begin{proof}
The classification of archimedean absolute values gives an exponent in
$(0,1]$. In characteristic zero, evaluation at $2$ gives
\[
 \kappa(\omega)=\frac{\max(0,\log\phi_\omega(2))}{\log2}.
\]
Indeed the archimedean value of $2$ is $2^{\kappa(\omega)}$, whereas a
nonarchimedean value of $2$ is at most one. The formula proves
measurability and integrability by the defining properties of an adelic
structure. Its strictly positive locus is precisely the archimedean locus,
which is consequently measurable. In positive characteristic there are no
archimedean places, and $\kappa$ is the zero function, giving all assertions.
\end{proof}
